\documentclass[../../main.tex]{subfiles}
\begin{document}

\begin{center}
  {\fontsize{48pt}{58pt}\selectfont \textbf{2013}\normalsize\textbf{年度}}
\end{center}

\begin{table}[htb]
  \centering
  \caption{出題テーマと難易度}
  \begin{tabular}{|c|p{9cm}|c|}
    \hline
    \textbf{問題番号} & \textbf{テーマ} & \textbf{難易度} \\
    \hline
    第1問 & 漸化式と整除，さいころの確率 & \\
    \hline
    第2問 & 行列式・跡と行列の五乗 & \\
    \hline
    第3問 & 指数関数を含む方程式の正の解の個数 & \\
    \hline
    第4問 & 三角関数の不等式を満たす区間の長さと極限 & \\
    \hline
    第5問 & 円と楕円の内接条件，接点と面積 & \\
    \hline
  \end{tabular}
\end{table}

\begin{table}[htb]
  \centering
  \caption{解答}
  \renewcommand{\arraystretch}{1.4}
  \begin{tabular}{|c|c|p{9cm}|}
    \hline
    \textbf{問題番号} & \textbf{小問} & \textbf{答え} \\
    \hline
    第1問 & (1) & $5$の倍数になる（証明問題） \\
    \cline{2-3}
    & (2) & 確率 $\dfrac{325}{648}$ \\
    \hline
    第2問 & (1) & $\Delta(AB)=\Delta(A)\Delta(B)$（証明問題） \\
    \cline{2-3}
    & (2) & $x=1$，$y=2,\dfrac{-1+\sqrt{5}}{2},\dfrac{-1-\sqrt{5}}{2}$ \\
    \hline
    第3問 & -- & $k<0$で0個，$k=0$で1個，$0<k\le1$で2個，
    $1<k<e-1$で3個，$k=e-1$で2個，$k>e-1$で1個 \\
    \hline
    第4問 & -- & $\displaystyle\lim_{n\to\infty}S_n=\dfrac{\pi}{8}$ \\
    \hline
    第5問 & (1) & $b^2<\dfrac{1}{2}$なら
    $a^2=b^2(1-b^2)$，$b^2\ge\dfrac{1}{2}$なら$a=\dfrac{1}{2}$ \\
    \cline{2-3}
    & (2) & $(p,q)=\left(\dfrac{\sqrt{2}}{2},\dfrac{\sqrt{6}}{6}\right)$ \\
    \cline{2-3}
    & (3) & 面積
    $\sqrt{3}\left(\dfrac{\pi}{12}-\dfrac{1}{18}\right)-\dfrac{\pi}{27}$ \\
    \hline
  \end{tabular}
\end{table}

\paragraph{試験形態}
5問．試験時間は要確認．

\end{document}
